\section{Dynamic Baselining}

\label{sec:baselines}
% ══════════════════════════════════════════════════════════════════════════════

\subsection{The Baseline Problem}

Carbon credits quantify the difference between actual outcomes and a
counterfactual baseline---what would have happened without the
intervention. Traditional approaches allow project developers to
construct their own baselines, creating a fundamental conflict of
interest. Higher baseline deforestation projections yield more credits,
incentivizing inflated baselines \citep{west2023action}.

\subsection{Synthetic Control Approach}

We adopt the synthetic control method \citep{abadie2010synthetic} for
baseline construction. For each project area $P$, we identify a set of
$J$ control areas $\{C_1, \ldots, C_J\}$ that matched the project area's
pre-intervention characteristics but were not subject to conservation
intervention.

The synthetic control baseline is:
\begin{equation}
  \hat{B}_t^{\text{baseline}}(P) = \sum_{j=1}^{J} w_j \cdot B_t(C_j)
\end{equation}
where weights $w_j \geq 0$, $\sum w_j = 1$ are chosen to minimize
pre-intervention prediction error:
\begin{equation}
  \mathbf{w}^* = \arg\min_{\mathbf{w}} \sum_{t < t_0}
  \left| B_t(P) - \sum_j w_j B_t(C_j) \right|^2
\end{equation}

\subsection{Control Area Selection}

Control areas are selected from a candidate pool using the following
matching criteria:
\begin{itemize}
  \item Same biome and eco-region (WWF classification).
  \item Similar deforestation pressure (road density, population growth,
    commodity prices).
  \item Similar forest type and initial biomass ($\pm$20\%).
  \item No conservation intervention during the study period.
  \item Minimum 50\,km buffer from the project boundary.
\end{itemize}

\subsection{Baseline Validation}

We validate dynamic baselines against actual deforestation in 23
projects where the true counterfactual is observable (projects that
lost funding and were abandoned). Dynamic baselines predict actual
deforestation with $R^2 = 0.87$, compared to $R^2 = 0.41$ for
project-developer baselines and $R^2 = 0.72$ for jurisdictional
averages.


% ══════════════════════════════════════════════════════════════════════════════
